% Literature Review. Rewritten 2026-10-04 from reports/Literature review quality assessment.md and
% revised the same day from reports/Literature review quality assessment 2.md (plan:
% ~/.claude/plans/virtual-herding-flute.md). Modelled on Jimenez Ordonez (2026) Sec. 1.1, see
% thesis/style_guide_sota_aida.md. Re-voiced 2026-10-04 (user): the phrasing and paragraph wrap-ups
% of the first version ("Using ... relies on", "Computational studies have also explored",
% "Machine learning has also been used", "Despite these advancements", "Bridging this gap is
% essential", "The integration of ...") with every correction and added study of the revision kept.
% Table: four studies relating a measured geometry to disease in patients (as Aida's four-row table
% covers one strand); columns make the gap visible (cohort, approach, timing).
% Condensed 2026-10-04 to fit 3 pages (user): material already in the introduction dropped (Asakura
% detail, Temov variation, SCORE2 description); Bekircavusoglu, Zebic Mihic, Givehchi and the
% Chatzizisis mechanism no longer cited here (Zebic Mihic and Temov remain in the introduction).
% Order: S1 history + Asakura (Korona pattern), van der Giessen, Stone 2012 | Kashyap, Zhang, Garcha (chronological) |
% Griffo | S2 Cui (choice of angle; Givehchi dropped 2026-10-04, measurement error is not the point) | tortuosity (Groves, Li, Zebic Mihic 2024, Nannini, Tello Ayala;
% reconciled by endpoint) | extraction (Bransby, Nannini), per-artery scope | S3 Temov variation,
% Juan, Bekircavusoglu, Tommasino | Tello Ayala transformer + Glagov | Han | Sakellarios, Won |
% SCORE2, Chan, Bergstrom | gap | search statement | table.
% shen2026geometry deliberately NOT cited anywhere in the thesis (user, 2026-10-04).
% Rules: author and year as subject, \cite after the year; no numbers in the prose (cohort sizes in
% the table only); no concept explanations; no inserted clauses between two commas; no em dashes;
% synthesis only when every element is cited.
% Full text read: kashyap, li, griffo, telloayala, zhang2024curvature, won2022paradigm, chan2024orfan,
% bergstrom2021scapis, han2022plaque, tommasino2024clap, cui2017bifurcation, zebicmihic2024index,
% nannini2024characterization, juan2017bifurcation, garcha2025sensitivity; vandergiessen and
% sakellarios via the publisher page (quotes to confirm by eye).
% Mechanism (Chatzizisis) lives in clinical_background_v2/03_wall_shear_stress.tex; Murray in
% methodology/08_daughter_ratio.tex; curvature and tortuosity index in technical_background/04_curvature.tex.
% 2026-10-07: Sec. 1 opened on an uncited Aida-style frame (depends on) plus field topic sentence; Stone 2012 restored as the prospective step after van der
% Giessen; Kashyap paragraph opened on a claim; synthesis no longer opens with "Collectively ... highlight".
% Abstract-level only: asakura1990flow, stone2012prediction, groves2009tortuosity,
% bekircavusoglu2026ramus, temov2016bifurcation, glagov1987remodeling, score2_2021.
% TODO (author): search statement needs the databases searched, the time window and search date, an
% inclusion criterion, and one reason per exclusion (see report 2, Priority 4).
% STILL OWED (user chose already-read sources only, 2026-10-04; read before adding):
%   near misses for the gap: Chen 2024 EHJ-CVI (doi 10.1093/ehjci/jeae135, paywalled), Lee 2024 JCCT
%   PARADIGM radiomics (doi 10.1016/j.jcct.2024.02.003), CArTI AHA 2025 abstract
%   (doi 10.1161/circ.152.suppl_3.4365988), SCAPIS re-examination (Good 2026, doi 10.1111/joim.70068);
%   WSS outcomes: Lee 2019 EMERALD, Stone 2018 PROSPECT ESS; calcium scoring: ESC 2021 or Detrano 2008;
%   artery labelling: Wu 2019 TreeLab-Net, hampe2024gnn; radius ratios: taylor2024murray.
% 2026-10-05: gap stated once (end of Sec. 3); closing gap sentences of Sec. 1 and 2, the Asakura
% mechanism retelling and the Sec. 3 opener that duplicated the introduction removed. Zebic Mihic
% no longer cited in the introduction either. No forward references (user).
% CAD and CCTA are defined in the introduction. Packages: booktabs, array (table).

\chapter{Literature Review}
\label{ch:literature}

\section{Coronary Geometry and Wall Shear Stress}
\label{sec:lit-wss}

Establishing coronary geometry as a risk factor for CAD fundamentally depends on understanding how
the shape of a vessel affects plaque formation in its wall.
Research on this link has largely focused on wall shear stress (WSS) as the pathway from vessel
geometry to plaque.
Early foundational work by Asakura and Karino (1990) \cite{asakura1990flow} demonstrated that
plaque concentrates at bifurcations and bends in autopsy hearts.
Van der Giessen et al.\ (2009) \cite{vandergiessen2009bifurcation} confirmed this pattern in living
patients by locating plaque on the outer walls of bifurcations, although they inferred low WSS from
the position of the plaque rather than computing it.
Stone et al.\ (2012) \cite{stone2012prediction} computed WSS in patients after an acute coronary
syndrome and showed that low WSS at the first scan predicted where plaque progressed and the lumen
narrowed at follow-up.

Simulations have since identified which features of the geometry produce low WSS.
Kashyap et al.\ (2022) \cite{kashyap2022tortuosity} reconstructed the left main bifurcation from the
CCTA of individuals without CAD and simulated the blood flow through each.
They then compared every tortuosity measure used in earlier studies with the share of the vessel
wall exposed to low WSS.
Average curvature was the measure most closely related to low WSS, and the relationship was
significant although weak.
The tortuosity index showed no relationship with low WSS.

Zhang et al.\ (2024) \cite{zhang2024curvature} extended such simulations from a single bifurcation to
whole left coronary trees reconstructed from a public CCTA dataset, comparing stenosis-free trees
with moderately and severely stenosed ones.
They measured geometry and WSS for the whole tree and separately for the segments at and between
bifurcations.
Average curvature separated stenosed from stenosis-free trees at the whole-tree level and between
bifurcations, but not at the bifurcations themselves.
The area of low WSS separated only moderate from severe stenosis.

Garcha and Grande Guti\'errez (2025) \cite{garcha2025sensitivity} instead generated a large set of
artificial left coronary trees, which let them change one feature at a time and simulate the flow
in each.
Between trees they varied the branch angle, the bending of the vessels and the division of diameter
between the left anterior descending (LAD) and left circumflex (LCX) arteries.
The area of low WSS at the bifurcation grew as the diameters of the two branches became more unequal.
An uneven division also strengthened the effect of the angle and the bending, so these features act
together rather than independently.

Griffo et al.\ (2026) \cite{griffo2026wss} showed that the shape of the vessel alone is enough to
estimate WSS, by training a graph neural network to do so without simulating the flow.
The network learned from single coronary arteries reconstructed from invasive angiograms, with
simulated WSS as the reference, and its estimate agreed closely with the simulation, more so once
each vessel's WSS was divided by its own average.
In a test set of patients who later had a heart attack, the estimated and the simulated WSS
identified the lesion that caused it equally well, although both did so only moderately.

Vessel shape therefore decides where WSS is low, and it carries enough of the flow information for
WSS to be estimated without simulating it
\cite{kashyap2022tortuosity,garcha2025sensitivity,zhang2024curvature,griffo2026wss}.

\section{Quantifying Coronary Geometry}
\label{sec:lit-measure}

Using coronary geometry as a marker relies on a consistent quantification of vessel shape.
However, clinical studies have quantified coronary geometry inconsistently over the years, and the
chosen measure can determine whether a relationship is detected at all.

Cui et al.\ (2017) \cite{cui2017bifurcation} measured three angles at the left main bifurcation on
the CCTA of patients referred for invasive angiography.
These were the angle between the LAD and LCX and the angles between the left main artery and each of
its two daughters.
Only the angle between the LAD and LCX was wider in patients with significant stenosis, while the two
angles involving the left main artery did not differ.
A study measuring only the angles involving the left main artery would therefore have found no
relationship between the bifurcation angle and stenosis.

For tortuosity, studies using different measures have disagreed even on the direction of the
association.
Groves et al.\ (2009) \cite{groves2009tortuosity} screened the invasive angiograms of consecutive
patients and graded tortuosity visually as two consecutive turns that reversed the direction of a
major artery.
Patients with such severe tortuosity were less likely to have significant CAD than a randomly chosen
group without it.
Li et al.\ (2011) \cite{li2011tortuosity} instead counted sharp bends along the main arteries of
consecutive patients with chest pain.
Tortuous patients were again less likely to have CAD, although the association was borderline, and
among patients with CAD the rate of adverse events at follow-up did not differ with tortuosity.
Most recently, Tello Ayala et al.\ (2026) \cite{telloayala2026tortuosity} segmented the right
coronary artery automatically in a large archive of angiograms from referred patients and computed
the turning angle at each point along its centerline.
Greater average tortuosity was associated with CAD, and a transformer that read the full profile of
turning angles identified CAD better than the average alone.
Depending on how tortuosity is measured, the same artery could therefore be read as a sign of less
disease or of more.

The geometric studies reviewed here examined only the right coronary artery or the left
tree alone \cite{telloayala2026tortuosity,garcha2025sensitivity,zhang2024curvature}.
To determine the full potential of the coronary tree as a risk marker, this must be extended.
Collectively, these findings suggest that an inconsistent measure
\cite{cui2017bifurcation,kashyap2022tortuosity} and a partial view of the tree may each hide where
geometry relates to disease.

Using geometry in routine risk assessment relies on automated segmentation of the coronary tree from
CCTA.
Bransby et al.\ (2026) \cite{bransby2026imagecasx} annotated a large public CCTA dataset anew and
benchmarked several deep learning segmentation methods on it.
They compared each method with the agreement between expert annotators of the same scans.
The best method came close to this agreement but ultimately concluded that automated
segmentation does not yet match human annotators.

\section{Geometry and Risk of Coronary Disease}
\label{sec:lit-risk}

Estimating an individual's risk of developing plaque from geometry requires geometry measured before
plaque appears and a later scan showing whether it does.
Geometry measured once plaque is present cannot show which came first, as an artery remodels while
plaque grows \cite{glagov1987remodeling}.

Studies with follow-up have mostly concerned disease that was already present.
Sakellarios et al.\ (2017) \cite{sakellarios2017bifurcation} followed bifurcations with repeated
CCTA in a small group of patients with established disease, reporting that WSS predicted lumen
reduction best when the daughter branch was included.
Han et al.\ (2022) \cite{han2022plaque} compared a small sample of patients who later had an acute
coronary syndrome with matched controls from a registry of patients referred for suspected disease.
Their findings indicate that the lesions that caused the event lay nearer the ostium and more often
at bifurcations than other lesions in the same patients.
However, the typical event followed within weeks of the scan, so these lesions were already present
at imaging.
Tommasino et al.\ (2024) \cite{tommasino2024clap} followed patients referred for CCTA for one year,
linking a wide left main bifurcation angle to progression of proximal LAD disease in a cohort in which
many already had proximal LAD disease.
Most of these events were progression on imaging requested by clinicians, and their reported
intervals cannot be reproduced.
In contrast, Won et al.\ (2022) \cite{won2022paradigm} followed registry patients without plaque on
a first CCTA to a repeat scan some years later, observing new plaque in a substantial share of them
and relating plaque progression to age, sex, hypertension, diabetes, hyperlipidemia, obesity,
smoking and HbA1c but not to geometry.

Despite these advancements, a significant gap remains in the current literature.
Cohorts suited to this question have not measured geometry, since the plaque-free registry of Won
et al.\ included none and the random population sample of Bergstr\"om et al.\
\cite{bergstrom2021scapis} was scanned once and so had no follow-up.
No study identified in this search has combined automated measurement of the whole coronary tree with
follow-up to new plaque in a population selected at random rather than referred for suspected
disease.
An overview of the studies relating coronary geometry to disease is given in
Table~\ref{tab:literature}.

\begin{table}[htbp]
  \centering
  \footnotesize
  \begin{tabular}{@{}>{\raggedright\arraybackslash}p{0.16\linewidth}%
                  >{\raggedright\arraybackslash}p{0.22\linewidth}%
                  >{\raggedright\arraybackslash}p{0.25\linewidth}%
                  >{\raggedright\arraybackslash}p{0.29\linewidth}@{}}
    \toprule
    Author \& Year & Cohort & Geometric Approach & Main Findings \\
    \midrule
    Zhang (2024) & 39 left coronary trees, stenosed and stenosis-free
      & Whole-tree curvature on CCTA, cross-sectional
      & Curvature separates stenosed from stenosis-free trees. \\
    Li (2011) & 1010 consecutive angiography patients
      & Bend count on invasive angiograms, cross-sectional
      & Tortuosity inversely associated with CAD. \\
    Tello Ayala (2026) & 22,334 referred patients
      & Turning-angle profile along the RCA, cross-sectional
      & Tortuosity higher with CAD; the full profile identifies CAD better than its mean. \\
    Han (2022) & 116 referred patients with later ACS, 116 controls
      & Lesion position and shape on CCTA, follow-up to ACS
      & Event-causing lesions nearer the ostium and at bifurcations; events within weeks. \\
    \bottomrule
  \end{tabular}
  \caption{Overview of studies relating coronary geometry to disease, highlighting the cohort, the
    geometric approach and the main findings. None measured geometry in a randomly selected
    population and followed it to new plaque. ACS: acute coronary syndrome; CAD: coronary artery
    disease; CCTA: coronary computed tomography angiography; RCA: right coronary artery.}
  \label{tab:literature}
\end{table}
